\section*{Methods}

\paragraph{Platform and ground-truth parameters.} Illustrative calculations use the CN Bio PhysioMimix hardware parameters reported by Docci et al.\ \citep{docci2022}: media volume $V_m = 1600\ \mu$L, cell-compartment volume $V_c = 3\ \mu$L (a nominal small volume, not independently reported), $3\times10^5$ cells. Ground-truth kinetic parameters ($CL_d = 8.0\ \mu$L/min, $K_p = 5.0$, $k_b = 0.0008$ min$^{-1}$) are illustrative, chosen to be broadly consistent with the orders of magnitude reported for liver-chip systems, and are not fitted to any specific compound. Two $CL_{int}$ levels were used: a ``low'' level (total $CL_{int} = 0.3\ \mu$L/min, corresponding to $1.0\ \mu$L/min per $10^6$ cells) and a ``high'' level (total $CL_{int} = 15.0\ \mu$L/min, corresponding to $50.0\ \mu$L/min per $10^6$ cells), spanning part of the $\sim$800-fold range of intrinsic clearances reported across compounds by Rajan et al.\ \citep{rajan2023}.

\paragraph{Structural identifiability and the equivalence family.} Equation~\eqref{eq:threeeqns} was solved symbolically (SymPy) for $(CL_d, K_p, CL_{int})$ in terms of a free $k_b$, yielding Eq.~\eqref{eq:family}. This closed-form family was verified numerically by substituting two specific parameter quadruples differing up to 4-fold --- the ground-truth quadruple and a second quadruple obtained from Eq.~\eqref{eq:family} at a different $k_b$ --- into the full nonlinear ODE system (Eqs.~\eqref{eq:cm}--\eqref{eq:cc}, constant volume, $e=0$) and numerically integrating both (\texttt{scipy.integrate.solve\_ivp}, method LSODA, relative tolerance $10^{-9}$, absolute tolerance $10^{-12}$) over a simulated 2000-minute incubation. The two resulting $C_m(t)$ trajectories were compared directly (Fig.~\ref{fig:degeneracy}).

\paragraph{Exact asymptotic limits.} Equations~\eqref{eq:fastdist}--\eqref{eq:fastmetab} were derived analytically from the eigenvalues of $A$ in the limits $CL_d\to\infty$ and $CL_{int}\to\infty$ respectively, and confirmed by direct numerical eigenvalue evaluation across $CL_d \in \{8, 80, \dots, 8\times10^{12}\}\ \mu$L/min and $CL_{int} \in \{4.5, 45, \dots, 4.5\times10^{10}\}\ \mu$L/min (other parameters fixed at ground truth), matching the closed-form limits to 5 significant figures before floating-point noise appeared. The nondimensionalization (Eq.~\eqref{eq:nondim}) was derived by substitution and confirmed symbolically.

\paragraph{Synthetic data generation.} For a given platform, drug, dose, and sampling schedule, noise-free $C_m(t)$ was simulated by numerical integration as above. Proportional Gaussian noise (coefficient of variation 15\%, representative of LC-MS/MS assay variability) was added to generate synthetic ``observed'' data, with a fixed random seed for reproducibility. Sparse designs used 5 samples and dense designs used 12 samples, log-spaced or evenly spaced across an incubation window chosen adaptively per scenario (shorter for high-clearance compounds, consistent with the use of shorter incubation windows for high-clearance compounds and longer windows for low-clearance compounds in practice \citep{rajan2023}).

\paragraph{Mechanistic model fitting.} All five kinetic/loss parameters ($CL_{int}$, $CL_d$, $K_p$, $k_b$, evaporation rate $e$) were fitted jointly to synthetic data by nonlinear least squares on log-transformed residuals (\texttt{scipy.optimize.least\_squares}, trust-region-reflective algorithm, bounded), starting from perturbed values near the ground truth as a literature-informed analyst's initial guess would be.

\paragraph{Practical identifiability: Fisher information.} At the fitted optimum, the Gauss-Newton Fisher Information Matrix $J^TJ$ was computed from the Jacobian of the residuals. The relative standard error of $CL_{int}$ was computed two ways: via plain matrix inversion (\texttt{numpy.linalg.inv}) and via a truncated-SVD pseudoinverse (\texttt{numpy.linalg.pinv}, \texttt{rcond}=$10^{-10}$), for two representative scenarios (low- and high-$CL_{int}$, dense sampling). The condition number of $J^TJ$ was computed for both scenarios (Fig.~\ref{fig:fim}).

\paragraph{Practical identifiability: profile likelihood.} For the low-$CL_{int}$, sparse-sampling scenario, a profile likelihood for $CL_{int}$ was computed by fixing $CL_{int}$ at a grid of values relative to its fitted value (ratios 0.2 to 19.2), re-optimizing the remaining free parameters at each grid point from a single fixed starting point (\emph{single-start}), and recording the resulting sum of squared residuals (SSE) relative to the unconstrained optimum. This was repeated with 6 random restarts per grid point, retaining the best (lowest-SSE) fit at each point (\emph{multi-start}), for three scenarios: sparse sampling with $k_b$ and evaporation $e$ unknown (jointly fitted); sparse sampling with $k_b$ and $e$ fixed at their true values; and dense sampling with $k_b$ and $e$ fixed. A 95\% confidence region was defined by the $\chi^2_1$ threshold ($\Delta$SSE $\le 3.84\,\sigma^2$, with $\sigma^2$ the residual variance at the optimum) (Fig.~\ref{fig:profile}).

\paragraph{Software and reproducibility.} All calculations were performed in Python 3 (NumPy, SciPy, SymPy, Matplotlib). Code implementing every result in this paper, including the figures, is provided as Supplementary Material and will be deposited in a public repository.
